\section{Finite-field computations without known roots}
\label{ff:section}

The splitting procedures will manipulate all roots of a polynomial at once.
Two elementary tools make this possible.  First, a computation over a field
can be executed in a product of fields until a zero test distinguishes two
components.  Second, a supplied generator of a primary multiplicative subgroup
permits logarithms and root extraction by short digit calculations.  We develop
both tools explicitly, and then use them to split every totally split polynomial
of even degree.

Finite fields are represented by a basis over their prime field and a
multiplication table, or by an equivalent tower of polynomial quotients.
Field arithmetic, binary powering, polynomial Euclidean algorithms, and linear
algebra are therefore deterministic procedures with polynomial bit complexity
in the representation lengths.  All basis choices and pivot choices below use
a fixed ordering.  For a prime $q$ and a positive integer $m$, write
$v_q(m)$ for its $q$-adic valuation.  If $k$ has cardinality $Q$ and
$q\mid Q-1$, its full $q$-primary subgroup is the unique subgroup of
$k^*$ of order $q^{v_q(Q-1)}$.

\subsection{Simultaneous execution and zero tests}

Let $k$ be an explicitly represented finite extension of $\mathbf F_p$, and let
$F\in\mathbf F_p[T]$ be monic, square-free, and totally split, of degree $N$.
Write its roots as $x_1,\ldots,x_N$ for the proof; the algorithm does not know
these roots.  The algebra
\begin{equation}
 A=k[T]/(F)\simeq\prod_{i=1}^N k,
 \qquad a(T)\longmapsto(a(x_1),\ldots,a(x_N))
 \label{ff:product}
\end{equation}
lets us carry out arithmetic on all components simultaneously.

\begin{lemma}[Simultaneous execution]
\label{ff:simultaneous}
Suppose a deterministic field procedure uses arithmetic, zero tests, and
integer control, and its scalar input at component $i$ is obtained by
specializing given elements of $A$ at $T=x_i$.  Assume the procedure is
defined for each specialized input.  One can either execute the
procedure on all components with a common sequence of decisions, or return a
nontrivial proper monic divisor of $F$ in $\mathbf F_p[T]$.
The extra cost of each scalar operation or zero test is polynomial in
$N$, $[k:\mathbf F_p]$, and $\log p$.
\end{lemma}

\begin{proof}
For a zero test on $a\in A$, compute $d=\gcd(F,a)$ over $k$ using its
representative of degree less than $N$.  If $d=1$, every component is nonzero;
if $d=F$, every component is zero.  In the remaining case, $d$ is a proper
factor and the procedure stops.  Although computed over $k$, its monic form
belongs to $\mathbf F_p[T]$: it is the product of $T-x_i$ over the components
where $a(x_i)=0$.

Addition and multiplication take place in $A$.  A nonzero scalar is inverted
only after the zero test has shown it nonzero in every component; the extended
Euclidean algorithm then gives its inverse modulo $F$.  Thus, until a split
occurs, every field instruction and decision is valid in every component.
Polynomial degree tests and choices of pivots reduce to successive zero tests
on coefficients, so the same argument applies inside polynomial arithmetic
and linear algebra.  The usual dense implementations give the cost bound.
\end{proof}

This is the dynamic-evaluation principle~\cite{DellaDoraDicrescenzoDuval85};
we call its use here \emph{simultaneous execution}.  In particular, a
divergent pivot or a divergent polynomial degree is useful information: it
already supplies the required factor.  The convention also applies to nested
procedures whose coefficient field is being simulated by $A$.

\subsection{Logarithms in primary subgroups}

Let $q$ be a prime, and suppose $u$ generates a cyclic group of order $q^s$.
For $a\in\langle u\rangle$, its logarithm $\ell\in[0,q^s)$ can be recovered
one base-$q$ digit at a time, as in the prime-power logarithm algorithm of
Pohlig and Hellman~\cite{PohligHellman78}.  If $\ell_j\equiv\ell\pmod{q^j}$ is already
known, the next digit is the unique $c\in\{0,\ldots,q-1\}$ satisfying
\begin{equation}
 (au^{-\ell_j})^{q^{s-j-1}}
   =(u^{q^{s-1}})^c,
 \qquad 0\le j<s.
 \label{ff:digit}
\end{equation}
Set $\ell_{j+1}=\ell_j+cq^j$ and begin with $\ell_0=0$.
Indeed, if $\ell-\ell_j=q^j(c+qv)$, raising to $q^{s-j-1}$ removes the term
$qv$ and leaves exactly Equation~\eqref{ff:digit}.  This uses $q$ equality
tests per digit, rather than an enumeration of $q^s$ group elements.

The same calculation works in a product of fields when $u$ has order $q^s$
in every component.  When the product is presented as $k[X]/H$ with $H$
square-free, gcds with $H$ partition the components at each digit according
to the equality in Equation~\eqref{ff:digit}; retain each nonempty piece and
its current digit string.  There are at most $\deg H$ pieces at every level.
Keep each final integer logarithm together with its polynomial modulus.
Algebra elements subsequently computed from these logarithms can be recombined
by polynomial Chinese remaindering.  Alternatively, when zero components are represented by idempotent masks,
perform the same partitions using those masks.  These two implementations
will be used below.

\subsection{Root extraction in a finite reduced algebra}

We next give the root procedure used in the norm equations later in the paper.
Its input already includes a generator over the constant field, but it does
not include generators in extension fields, nor a factorization of the
modulus.

\begin{lemma}[Root extraction from a primary generator]
\label{ff:unitroot}
Let $k$ be an explicit finite field of cardinality $Q$ and characteristic $p$.
Let $q$ be an odd prime dividing $Q-1$, let $\zeta\in k$ have order $q$, and
let $\omega\in k$ generate the full $q$-primary subgroup of $k^*$.
Let $H\in k[X]$ be monic and square-free, with $h=\deg H\ge1$, and suppose
$p>h^2$.  Given a unit $a\in k[X]/H$ that is promised to be a $q$-th power,
one can deterministically construct $b$ with $b^q=a$.
The bit cost is polynomial in $q$, $h$, and the length of the given field
representation.  The procedure uses only arithmetic and zero tests over $k$,
and so admits simultaneous execution as in Lemma~\ref{ff:simultaneous}.
\end{lemma}

\begin{proof}
The construction has three steps: separate extension degrees, obtain a primary
generator on each resulting part, and take the two coprime parts of the root.

\smallskip
\noindent\emph{Separate extension degrees.}
For $d=1,\ldots,h$, use gcds with $X^{Q^d}-X$, removing factors already found,
to form the product $H_d$ of the irreducible factors of $H$ of degree $d$.
All powers are computed modulo the current modulus by binary powering.
Discard $H_d=1$.  Every field component of $k[X]/H_d$ has cardinality $Q^d$.
Write
\[
 Q^d-1=q^s t,\qquad q\nmid t.
\]
Only division by the specified prime $q$ is needed to find $s$ and $t$.

\smallskip
\noindent\emph{Obtain a primary generator.}
Since $q$ is odd and $q\mid Q-1$, the elementary lifting-the-exponent identity
is
\begin{equation}
 v_q(Q^d-1)=v_q(Q-1)+v_q(d).
 \label{ff:lte}
\end{equation}
To see this, raising a number congruent to $1$ modulo $q$ to an exponent
prime to $q$ preserves the valuation of its difference from $1$, whereas
raising it to the $q$-th power increases that valuation by exactly one.
Both assertions follow from the binomial expansion: its linear term has
strictly smaller $q$-valuation than all subsequent terms, using that $q$
is odd in the second assertion.  Iterating proves Equation~\eqref{ff:lte}.
If $q\nmid d$, set $u=\omega$, which already has order $q^s$ in every
component.  Suppose instead that $q\mid d$.  Solve the $k$-linear equation
\begin{equation}
 v^{Q^{d/q}}=\zeta v
 \quad\text{in }k[X]/H_d.
 \label{ff:eigenvector}
\end{equation}
In each component its nonzero solutions exist: the cyclic group of order
$Q^d-1$ contains a solution of $v^{Q^{d/q}-1}=\zeta$, because
\[
 \zeta^{(Q^d-1)/(Q^{d/q}-1)}=1.
\]
For such a solution, Equation~\eqref{ff:lte} shows that the $q$-valuation of
its order is $s$.

To choose a solution nonzero in every component without knowing those
components, take a kernel basis $v_0,\ldots,v_{r-1}$ and try
\begin{equation}
 v(c)=\sum_{j=0}^{r-1}c^jv_j,
 \qquad c=0,\ldots,h^2.
 \label{ff:unit-search}
\end{equation}
Test each candidate for being a unit by a gcd with $H_d$.  In any component,
at least one basis vector is nonzero, so at most $r-1$ values of $c$ make
$v(c)$ zero there.  There are at most $h$ components and $r\le h$; fewer than
$h^2$ values are excluded altogether.  The hypothesis $p>h^2$ makes the
listed parameters distinct.  Hence the search succeeds.  Raising its result
to $t$ produces a generator $u$ of order $q^s$ in every component.

\smallskip
\noindent\emph{Take the root.}
Use the integer Chinese remainder theorem for $q^s$ and $t$ to choose
exponents projecting $a$ to its primary and prime-to-$q$ parts.  Explicitly,
put
\[
 a_q=a^{t(t^{-1}\bmod q^s)},\qquad
 a_t=a^{q^s((q^s)^{-1}\bmod t)},
\]
with $a_t=1$ when $t=1$.  Then $a=a_q a_t$, the order of $a_q$ divides
$q^s$, and the order of $a_t$ divides $t$.  A root of $a_t$ is obtained by
raising it to $q^{-1}\bmod t$, again using $1$ when $t=1$.

Compute the logarithm of $a_q$ to $u$ by Equation~\eqref{ff:digit},
partitioning $H_d$ as necessary.  In each final piece its logarithm
$\ell\in[0,q^s)$ is divisible by $q$, because $a$ is a $q$-th power.
Therefore $u^{\ell/q}$ is a root of $a_q$ on that piece.  Multiply the two
roots, recombine the pieces by polynomial Chinese remaindering, and finally
recombine the relatively prime $H_d$.

For complexity, $d\le h$, the bit length of $Q^d$ is at most $1+h\log_2Q$,
and $s\le h\log_2Q$.  There are at most $h$ nonempty pieces at every digit
level; each digit tests at most $q$ possibilities.  Equation~\eqref{ff:eigenvector}
is a linear system of dimension at most $h$, and
Equation~\eqref{ff:unit-search} uses at most $h^2+1$ trials.  All remaining
operations are polynomial arithmetic, powering, or integer Euclidean
algorithms of the indicated sizes.  This proves the claimed polynomial
bound and also the assertion about field instructions.
\end{proof}

We denote this procedure by $\mathsf{UnitRoot}$.  To extract a root in $k$
itself, apply it with the linear modulus $H=X$.  The proof explains why
distinct-degree decomposition suffices: every subsequent partition is obtained
by an equality test, and full irreducible factorization is never invoked.

\subsection{Splitting an even number of roots}

R\'onyai proved that a supplied quadratic nonresidue suffices to split
a completely split polynomial of even degree in polynomial time~\cite{Ronyai88}.
A quadratic nonresidue gives a full $2$-primary generator by raising it to
the odd part of $p-1$.  We give the pair-orientation argument in the form
needed here.  A primary generator forces a distinction among an even number of
roots.  The mechanism is an orientation of the pairs of distinct roots: each
pair contributes one to precisely one of its two row counts.  Equal row
counts would then have the nonintegral value $(N-1)/2$.

\begin{proposition}[Even-degree splitting]
\label{ff:even-split}
Let $p$ be odd, and let $F\in\mathbf F_p[T]$ be monic, square-free, and
totally split of even degree $N\ge2$, with $p>N$.  Given a generator $u$ of
the full $2$-primary subgroup of $\mathbf F_p^*$, one can deterministically
find a nontrivial proper monic divisor of $F$ in a number of bit operations
polynomial in $N$ and $\log p$.
\end{proposition}

\begin{proof}
Set $A=\mathbf F_p[T]/F$ and $R=A\otimes_{\mathbf F_p}A$.  Write
$p-1=2^s t$ with $t$ odd.  In the coordinate description indexed by ordered
pairs of roots, the element
\[
 \Delta=T\otimes1-1\otimes T
\]
has coordinate $x_i-x_j$.  Thus $e=\Delta^{p-1}$ is the idempotent equal to
$1$ off the diagonal and $0$ on the diagonal.  Work in the algebra $eR$,
whose identity is $e$.  The element $\Delta^t$ has $2$-power order at every
coordinate there, so compute its logarithm to $u$ by the digit procedure.
For a masked equality test on an element $z$ in a subalgebra with identity
mask $e_0$, its zero mask is $e_0-(e_0z)^{p-1}$.  Hence every partition is
computed within $R$, with no knowledge of the roots.

Let $e_*$ be the sum of the final masks whose logarithms lie in
$[0,2^{s-1})$.  Replacing $(i,j)$ by $(j,i)$ changes $\Delta^t$ to its
negative, because $t$ is odd.  The logarithm therefore changes by
$2^{s-1}$ modulo $2^s$, and exactly one orientation of each unordered pair
is selected by $e_*$.  Apply the ordinary algebra trace on the second factor:
\begin{equation}
 c=(\operatorname{id}\otimes\operatorname{Tr}_{A/\mathbf F_p})(e_*)\in A.
 \label{ff:row-counts}
\end{equation}
Its $i$-th coordinate is the number $c_i$ of selected pairs in row $i$,
viewed in $\mathbf F_p$.  As integers,
\[
 0\le c_i\le N-1,
 \qquad \sum_{i=1}^N c_i=\frac{N(N-1)}2.
\]
If these integers were all equal, their common value would be $(N-1)/2$,
which is impossible for even $N$.  Since $p>N$, their images in the field
are not all equal either.  Testing $\gcd(F,c-a)$ for $a=0,\ldots,N-1$
therefore returns a nontrivial proper factor.

The algebra $R$ has dimension $N^2$, and the number of nonempty masks at a
digit level is at most $N^2$.  There are $s\le\log_2p$ levels, each with
two possible digits.  Binary powering, the trace computation, and the final
gcds all have the asserted polynomial cost.
\end{proof}
